\documentclass[10pt, aspectratio=169]{beamer}
\usepackage{../beamer_theme_408}

\title{408考研零基础 --- 操作系统}
\subtitle{3-1 操作系统是什么}
\author{}
\date{}

\begin{document}


\begin{frame}{本节目标}
    \begin{enumerate}
        \item 理解操作系统的定义与定位
        \item 掌握操作系统的四大功能
        \item 掌握操作系统的四大特征
        \item 了解操作系统的发展历程
    \end{enumerate}
\end{frame}

\begin{frame}{为什么需要操作系统}
    \begin{block}{问题：裸机为什么难以使用？}
        \begin{itemize}
            \item 用户直接操作硬件需要了解复杂的底层细节
            \item 多个程序同时运行时，资源分配与冲突管理复杂
            \item 缺乏统一的接口标准，软件移植困难
        \end{itemize}
    \end{block}
    \vspace{0.3em}
    \textbf{解决方案：引入中间层 --- 操作系统}
    \begin{center}
    \begin{tikzpicture}
        \draw[fill=blue!10,rounded corners] (0,0) rectangle (6,0.7);
        \node at (3,0.35) {\textbf{用户应用程序}};
        \draw[fill=orange!10,rounded corners] (0,-1) rectangle (6,-0.3);
        \node at (3,-0.65) {\textbf{\textcolor{408orange}{操作系统}}};
        \draw[fill=green!10,rounded corners] (0,-2) rectangle (6,-1.3);
        \node at (3,-1.65) {\textbf{计算机硬件（CPU、内存、磁盘...）}};
        \draw[<->,thick] (3,-0.3) -- (3,-1.3);
        \node[right,font=\small] at (6.2,-0.8) {系统调用/驱动};
    \end{tikzpicture}
    \end{center}
    \vspace{0.3em}
    \begin{itemize}
        \item OS隐藏硬件细节，提供统一的抽象接口
        \item 管理资源分配，确保多程序有序运行
    \end{itemize}
\end{frame}

\begin{frame}{操作系统定义}
    \textbf{操作系统（Operating System，OS）}是：
    \begin{itemize}
        \item 控制和管理计算机系统资源的\textbf{系统软件}
        \item 合理组织、调度计算机的工作与资源分配
        \item 为用户提供方便、有效的使用环境
    \end{itemize}
    \vspace{1em}
    \begin{block}{一句话定义}
        操作系统是介于\textcolor{408blue}{硬件}与\textcolor{408blue}{应用程序}之间的\textbf{系统软件}。
    \end{block}
    \begin{center}
        \begin{tabular}{c}
            \hline
            用户应用程序 \\
            \hline
            \textcolor{408orange}{操作系统} \\
            \hline
            计算机硬件 \\
            \hline
        \end{tabular}
    \end{center}
\end{frame}

\begin{frame}{操作系统的四大功能}
    \begin{enumerate}
        \item \textbf{处理机管理} —— CPU资源的分配与回收（进程/线程调度）
        \item \textbf{存储器管理} —— 内存的分配与回收、虚拟内存管理
        \item \textbf{设备管理} —— 管理I/O设备，完成用户的I/O请求
        \item \textbf{文件管理} —— 文件的存储、共享、保护与访问控制
    \end{enumerate}
    \vspace{0.5em}
    \begin{block}{目标}
        向上层用户和应用程序提供\textbf{方便、安全、高效}的服务。
    \end{block}
\end{frame}

\begin{frame}{四大功能与管理对象对应关系}
    \begin{table}
        \centering
        \begin{tabular}{|c|c|c|}
            \hline
            \textbf{功能} & \textbf{管理对象} & \textbf{核心问题} \\
            \hline
            处理机管理 & CPU（中央处理器） & 进程/线程调度、上下文切换 \\
            \hline
            存储器管理 & 内存（主存） & 地址转换、虚拟内存、页面置换 \\
            \hline
            设备管理 & I/O设备 & 缓冲区管理、SPOOLing、设备驱动 \\
            \hline
            文件管理 & 磁盘上的文件 & 目录结构、文件共享、访问控制 \\
            \hline
        \end{tabular}
    \end{table}
    \vspace{0.5em}
    \begin{block}{理解要点}
        每个功能对应一类硬件资源，OS的核心职责就是对\textbf{所有硬件资源进行抽象、分配、回收和保护}。
    \end{block}
\end{frame}

\begin{frame}{操作系统的四大特征}
    \begin{enumerate}
        \item \textbf{并发} —— 多个程序在\textbf{宏观上同时运行}，微观上交替执行
        \item \textbf{共享} —— 资源可被多个进程\textbf{互斥共享}或\textbf{同时共享}
        \item \textbf{虚拟} —— 将物理实体映射为多个逻辑对应物（虚拟处理器、虚拟内存）
        \item \textbf{异步} —— 进程的执行速度不可预知，但结果应一致
    \end{enumerate}
\end{frame}

\begin{frame}{并发与共享}
    \begin{columns}
        \begin{column}{0.48\textwidth}
            \textbf{并发}
            \begin{itemize}
                \item 单核CPU：交替执行（微观串行）
                \item 多核CPU：真正并行
                \item 与"并行"区别：\\
                      并发 $\neq$ 并行
            \end{itemize}
        \end{column}
        \begin{column}{0.48\textwidth}
            \textbf{共享}
            \begin{itemize}
                \item \textbf{互斥共享}：同一时刻只允许一个进程访问（如打印机）
                \item \textbf{同时共享}：同一时刻允许多个进程访问（如磁盘文件）
            \end{itemize}
        \end{column}
    \end{columns}
    \vspace{0.5em}
    \begin{alertblock}{注意}
        并发和共享是操作系统\textbf{最基本的两个特征}，二者互为存在条件。
    \end{alertblock}
\end{frame}

\begin{frame}{虚拟与异步}
    \textbf{虚拟}
    \begin{itemize}
        \item 虚拟处理器：通过分时复用，一个物理CPU虚拟为多个逻辑CPU
        \item 虚拟内存：通过请求调页，物理内存虚拟为更大的逻辑内存
        \item 虚拟设备：通过SPOOLing技术，独占设备虚拟为共享设备
    \end{itemize}
    \vspace{0.5em}
    \textbf{异步}
    \begin{itemize}
        \item 进程以不可预知的速度向前推进
        \item 只要资源充足，每次运行结果应相同（可再现性）
    \end{itemize}
\end{frame}

\begin{frame}{四大特征关系图}
    \begin{center}
    \begin{tikzpicture}[>=stealth]
        \node[draw=408blue,fill=408blue!10,rounded corners,minimum width=3cm,minimum height=1cm] (con) at (0,1.5) {\textbf{并发}};
        \node[draw=408blue,fill=408blue!10,rounded corners,minimum width=3cm,minimum height=1cm] (share) at (3.5,1.5) {\textbf{共享}};
        \node[draw=408orange,fill=408orange!10,rounded corners,minimum width=3cm,minimum height=1cm] (virt) at (0,0) {\textbf{虚拟}};
        \node[draw=408orange,fill=408orange!10,rounded corners,minimum width=3cm,minimum height=1cm] (async) at (3.5,0) {\textbf{异步}};
        \draw[<->,thick,408blue] (con) -- (share);
        \draw[->,thick,408blue] (con) -- (virt);
        \draw[->,thick,408blue] (share) -- (async);
        \node[below,font=\small] at (1.75,-0.5) {互为存在条件};
        \node[below,font=\small] at (1.75,-0.9) {并发和共享是最基本的两个特征};
    \end{tikzpicture}
    \end{center}
    \vspace{0.3em}
    \begin{itemize}
        \item \textbf{并发}是产生\textbf{共享}的前提，\textbf{共享}是\textbf{并发}的结果
        \item \textbf{虚拟}和\textbf{异步}是并发和共享的衍生特征
    \end{itemize}
\end{frame}

\begin{frame}{例题：判断OS特征}
    \textbf{例题1}：在一台单核CPU计算机上，用户同时打开了浏览器和Word，两个程序都在"运行"。请判断这体现了OS的什么特征？
    \vspace{0.5em}
    \begin{block}{解析}
        体现\textbf{并发}特征。单核CPU在某一时刻只能执行一个程序，浏览器和Word在宏观上同时运行，微观上交替使用CPU。
    \end{block}
    \vspace{0.5em}
    \textbf{例题2}：打印机在任何时刻只能为一个进程服务，这体现了OS的什么特征？
    \vspace{0.5em}
    \begin{block}{解析}
        体现\textbf{共享}特征中的\textbf{互斥共享}方式。打印机是临界资源，一次只能被一个进程访问，其他进程需要等待。
    \end{block}
\end{frame}

\begin{frame}{课堂练习}
    \textbf{判断以下场景体现的特征或功能：}
    \vspace{0.5em}
    \begin{enumerate}
        \item \textbf{场景A}：你的电脑同时播放音乐和编辑文档，感觉两个程序在"同时运行"。\\
        \textcolor{white!30!black}{\tiny 答案：并发}
        \vspace{0.3em}
        \item \textbf{场景B}：你的程序使用了4GB内存，但物理内存只有2GB，程序依然正常运行。\\
        \textcolor{white!30!black}{\tiny 答案：虚拟（虚拟内存/空分复用）}
        \vspace{0.3em}
        \item \textbf{场景C}：两个不同的程序同时读取同一个磁盘文件。\\
        \textcolor{white!30!black}{\tiny 答案：共享（同时共享）}
        \vspace{0.3em}
        \item \textbf{场景D}：操作系统从就绪队列中选择一个进程分配CPU。\\
        \textcolor{white!30!black}{\tiny 答案：处理机管理（进程调度）}
    \end{enumerate}
    \vspace{0.5em}
    \begin{alertblock}{提示}
        先自行思考，再对比答案。重点理解各特征的核心区分点。
    \end{alertblock}
\end{frame}

\begin{frame}{操作系统的发展历程}
    \begin{enumerate}
        \item \textbf{手工操作阶段} —— 用户独占全机，人机速度矛盾
        \item \textbf{单道批处理系统} —— 一次一道程序，串行执行，CPU利用率低
        \item \textbf{多道批处理系统} —— 多道程序并发，CPU利用率大幅提升
        \item \textbf{分时操作系统} —— 时间片轮转，交互性强
        \item \textbf{实时操作系统} —— 在规定时间内完成任务，硬实时/软实时
        \item \textbf{分布式操作系统} —— 多机协同，资源共享
    \end{enumerate}
\end{frame}

\begin{frame}{各种操作系统的比较}
    \begin{tabular}{|c|c|c|c|}
        \hline
        类型 & 主要目标 & 交互性 & 适用场景 \\
        \hline
        单道批处理 & 计算效率 & 弱 & 早期计算机 \\
        多道批处理 & CPU利用率 & 弱 & 大型计算 \\
        分时系统 & 响应时间 & 强 & 通用系统 \\
        实时系统 & 截止时间 & 中 & 工业控制 \\
        分布式系统 & 资源整合 & 强 & 网络环境 \\
        \hline
    \end{tabular}
\end{frame}

\begin{frame}{本节总结}
    \begin{enumerate}
        \item OS是控制和管理计算机资源的系统软件，提供方便高效的使用环境
        \item 四大功能：处理机管理、存储器管理、设备管理、文件管理
        \item 四大特征：并发、共享、虚拟、异步
        \item 发展历程：手工 $\to$ 单道批处理 $\to$ 多道批处理 $\to$ 分时 $\to$ 实时 $\to$ 分布式
    \end{enumerate}
\end{frame}

\begin{frame}{下节预告}
    \begin{center}
    \vspace{1em}
    {\Large \textbf{3-2 操作系统运行环境}}\\
    \vspace{0.3em}
    \textcolor{408blue}{了解操作系统运行的硬件基础}
    \vspace{1em}
    \end{center}
    \begin{block}{下节内容}
        \begin{itemize}
            \item \textbf{处理器状态}：内核态与用户态的区别与切换
            \item \textbf{中断机制}：外中断、内中断、异常的分类
            \item \textbf{系统调用}：用户程序请求OS服务的唯一入口
            \item \textbf{系统体系结构}：大内核与微内核的对比
        \end{itemize}
    \end{block}
    \vspace{0.5em}
    \begin{center}
        \textcolor{408orange}{敬请期待！}
    \end{center}
\end{frame}
\end{document}
